% !TEX program = xelatex
\documentclass[serif, aspectratio=169]{beamer}

\usepackage{fontspec}
\usepackage{polyglossia}
\setdefaultlanguage{vietnamese}
\setotherlanguage{english}

\newfontface\myfont{Times New Roman}

\usepackage{hyperref}
\usepackage{graphicx}
\usepackage{amsmath}
\usepackage{multicol}
\usepackage{booktabs}
\usepackage{lipsum}
\usepackage[sort, round]{natbib}
\bibliographystyle{plainnat}

\usepackage{graphicx,pstricks,listings,stackengine}

\author{Vũ Công Tuấn Dương -- B22DCKH024}
\title{Phân tích chuỗi thời gian:\\So sánh mô hình ARIMA và VAR}
\institute{
    Khoa CNTT1 -- PTIT
}
\date{\small \today}
\usepackage{HKUSTstyle}

\def\cmd#1{\texttt{\color{red}\footnotesize $\backslash$#1}}
\def\env#1{\texttt{\color{blue}\footnotesize #1}}

\definecolor{hkustyellow}{RGB}{167, 131, 55}
\definecolor{hkustblue}{RGB}{0, 56, 116}
\definecolor{hkustred}{RGB}{209, 51, 59}
\definecolor{codegreen}{RGB}{0, 128, 0}
\definecolor{codegray}{RGB}{128, 128, 128}
\definecolor{codepurple}{RGB}{128, 0, 128}

\lstset{
    basicstyle=\ttfamily\footnotesize,
    keywordstyle=\bfseries\color{hkustblue},
    commentstyle=\color{codegreen},
    stringstyle=\color{codepurple},
    numbers=left,
    numberstyle=\small\color{codegray},
    rulesepcolor=\color{red!20!green!20!blue!20},
    frame=single,
    breaklines=true,
    language=Python,
}

% ============================================================

\begin{document}

% --- Title Slide ---
\begin{frame}
    \titlepage
\end{frame}

% --- Table of Contents ---
\begin{frame}{Mục lục}
\tableofcontents[sectionstyle=show,
subsectionstyle=show/shaded/hide,
subsubsectionstyle=show/shaded/hide]
\end{frame}

% ============================================================
\section{Giới thiệu}
% ============================================================

\begin{frame}{Giới thiệu}
\begin{itemize}
    \item \textbf{Mục tiêu}: So sánh hai mô hình dự báo chuỗi thời gian:
    \begin{itemize}
        \item \textbf{ARIMA} -- Mô hình đơn biến (Univariate)
        \item \textbf{VAR} -- Mô hình đa biến (Multivariate)
    \end{itemize}
    \item \textbf{Bộ dữ liệu}: UCI Air Quality
    \begin{itemize}
        \item Nồng độ chất ô nhiễm theo giờ từ cảm biến tại Ý
        \item Thời gian: Tháng 3/2004 -- Tháng 4/2005
    \end{itemize}
    \item \textbf{Nguồn}: \url{https://archive.ics.uci.edu/dataset/360/air+quality}
\end{itemize}
\end{frame}

% ============================================================
\section{Tải và chuẩn bị dữ liệu}
% ============================================================

\subsection{Import thư viện}
\begin{frame}{Import thư viện}
\begin{lstlisting}
import numpy as np
import pandas as pd
import matplotlib.pyplot as plt
from statsmodels.tsa.arima.model import ARIMA
from statsmodels.tsa.vector_ar.var_model import VAR
from statsmodels.tsa.stattools import adfuller
from sklearn.metrics import mean_squared_error, mean_absolute_error
import warnings
warnings.filterwarnings('ignore')
\end{lstlisting}
\end{frame}

\subsection{Tải dữ liệu}
\begin{frame}{Tải dữ liệu}
\begin{lstlisting}
df_raw = pd.read_csv(
    'AirQualityUCI.csv',
    sep=';',
    decimal=',',
    na_values=-200,
)
df_raw = df_raw.dropna(how='all').reset_index(drop=True)

# Kết hợp Date + Time thành DatetimeIndex
df_raw['Datetime'] = pd.to_datetime(
    df_raw['Date'] + ' ' + df_raw['Time'].str.replace('.', ':'),
    format='%d/%m/%Y %H:%M:%S',
)
\end{lstlisting}
\end{frame}

\begin{frame}{Lựa chọn biến và resample}
Chọn 4 biến chính và lấy trung bình theo ngày:
\begin{lstlisting}
cols = ['CO(GT)', 'C6H6(GT)', 'NO2(GT)', 'T']
df = df_raw[['Datetime'] + cols].copy()
df = df.dropna(subset=cols).set_index('Datetime')
df = df.resample('D').mean().dropna()
\end{lstlisting}

\vspace{0.5cm}
\begin{table}
\centering
\begin{tabular}{ll}
\toprule
\textbf{Thông số} & \textbf{Giá trị} \\
\midrule
Kích thước & 341 ngày \\
Phạm vi ngày & 10/03/2004 $\rightarrow$ 04/04/2005 \\
Các biến & CO(GT), C6H6(GT), NO2(GT), T \\
\bottomrule
\end{tabular}
\end{table}
\end{frame}

% ============================================================
\section{Chia tập Train-Test}
% ============================================================

\begin{frame}{Chia tập Train-Test}
Tỷ lệ chia: \textbf{80\% Train -- 20\% Test}

\vspace{0.5cm}
\begin{table}
\centering
\begin{tabular}{lll}
\toprule
\textbf{Tập} & \textbf{Số ngày} & \textbf{Thời gian} \\
\midrule
Tổng cộng & 341 ngày & 10/03/2004 $\rightarrow$ 04/04/2005 \\
Train & 272 ngày & 10/03/2004 $\rightarrow$ 23/01/2005 \\
Test & 69 ngày & 24/01/2005 $\rightarrow$ 04/04/2005 \\
\bottomrule
\end{tabular}
\end{table}

\vspace{0.3cm}
\begin{lstlisting}
split_ratio = 0.8
split_idx = int(len(df) * split_ratio)
train = df.iloc[:split_idx]
test  = df.iloc[split_idx:]
\end{lstlisting}
\end{frame}

% ============================================================
\section{Kiểm định tính dừng (ADF)}
% ============================================================

\begin{frame}{Kiểm định tính dừng -- Augmented Dickey-Fuller}
\textbf{Kiểm định ADF}: Xác định chuỗi có dừng hay không.
\begin{itemize}
    \item $H_0$: Chuỗi có \textbf{đơn vị gốc} (không dừng)
    \item Nếu p-value $< 0.05$: \textbf{Bác bỏ} $H_0$ $\rightarrow$ Chuỗi dừng
\end{itemize}

\vspace{0.5cm}
\begin{table}
\centering
\begin{tabular}{lrrl}
\toprule
\textbf{Biến} & \textbf{p-value} & \textbf{Kết luận} \\
\midrule
CO(GT) & 0.0112 & \textcolor{codegreen}{\textbf{Dừng}} \\
C6H6(GT) & 0.0007 & \textcolor{codegreen}{\textbf{Dừng}} \\
NO2(GT) & 0.3339 & \textcolor{hkustred}{\textbf{Không dừng}} \\
T & 0.9046 & \textcolor{hkustred}{\textbf{Không dừng}} \\
\bottomrule
\end{tabular}
\end{table}

\vspace{0.3cm}
$\Rightarrow$ NO2 và Nhiệt độ cần sai phân ($d \geq 1$)
\end{frame}

% ============================================================
\section{Mô hình ARIMA}
% ============================================================

\subsection{Giới thiệu ARIMA}
\begin{frame}{Mô hình ARIMA}
\textbf{ARIMA(p, d, q)} -- Mô hình Tích phân Tự hồi quy Trung bình trượt:
\begin{itemize}
    \item \textbf{p}: Bậc thành phần Tự hồi quy (AR)
    \item \textbf{d}: Bậc sai phân (Integration)
    \item \textbf{q}: Bậc thành phần Trung bình trượt (MA)
\end{itemize}

\vspace{0.3cm}
\begin{block}{Đặc điểm}
\begin{itemize}
    \item Mô hình \textbf{đơn biến} -- mỗi chuỗi được mô hình hóa \textbf{độc lập}
    \item Không nắm bắt mối tương quan giữa các biến
    \item Tìm kiếm lưới (Grid Search) trên (p, d, q) để tối thiểu hóa AIC
\end{itemize}
\end{block}
\end{frame}

\subsection{Kết quả ARIMA}
\begin{frame}{Kết quả ARIMA -- Tìm tham số tối ưu}
\begin{table}
\centering
\begin{tabular}{lccc}
\toprule
\textbf{Biến} & \textbf{Bộ tham số $(p, d, q)$} & \textbf{AIC} \\
\midrule
CO(GT) & (3, 1, 3) & 619.14 \\
C6H6(GT) & (0, 1, 2) & 1532.69 \\
NO2(GT) & (2, 1, 3) & 2419.92 \\
T & (0, 1, 2) & 1258.95 \\
\bottomrule
\end{tabular}
\end{table}

\vspace{0.3cm}
\begin{lstlisting}
for col in cols:
    for p in [0, 1, 2, 3]:
        for d in [0, 1]:
            for q in [0, 1, 2, 3]:
                fitted = ARIMA(train[col], order=(p, d, q)).fit()
                # Chọn bộ tham số có AIC nhỏ nhất
\end{lstlisting}
\end{frame}

% ============================================================
\section{Mô hình VAR}
% ============================================================

\subsection{Giới thiệu VAR}
\begin{frame}{Mô hình VAR}
\textbf{VAR(k)} -- Mô hình Tự hồi quy Vector:
\begin{itemize}
    \item Mô hình \textbf{đa biến} -- tất cả biến được mô hình hóa \textbf{cùng lúc}
    \item Nắm bắt \textbf{tương quan chéo} giữa các biến
    \item Mỗi biến phụ thuộc vào giá trị quá khứ của \textbf{tất cả} biến
\end{itemize}

\vspace{0.3cm}
\begin{block}{Công thức}
\[
\mathbf{Y}_t = \mathbf{c} + \mathbf{A}_1 \mathbf{Y}_{t-1} + \mathbf{A}_2 \mathbf{Y}_{t-2} + \cdots + \mathbf{A}_k \mathbf{Y}_{t-k} + \boldsymbol{\varepsilon}_t
\]
với $\mathbf{Y}_t = [CO, C6H6, NO2, T]^T$
\end{block}
\end{frame}

\subsection{Kết quả VAR}
\begin{frame}{Kết quả VAR -- Chọn số bậc trễ tối ưu}
\begin{table}
\centering
\begin{tabular}{rr|rr|rr}
\toprule
\textbf{Bậc} & \textbf{AIC} & \textbf{Bậc} & \textbf{AIC} & \textbf{Bậc} & \textbf{AIC} \\
\midrule
1 & \textbf{7.18} & 6 & 7.40 & 11 & 7.55 \\
2 & 7.26 & 7 & 7.39 & 12 & 7.59 \\
3 & 7.23 & 8 & 7.42 & 13 & 7.68 \\
4 & 7.27 & 9 & 7.47 & 14 & 7.70 \\
5 & 7.34 & 10 & 7.51 & 15 & 7.77 \\
\bottomrule
\end{tabular}
\end{table}

\vspace{0.3cm}
\textbf{Kết luận}: Số bậc trễ tối ưu = \textbf{1} (AIC = 7.18)

\vspace{0.2cm}
\begin{lstlisting}
var_fitted = VAR(train).fit(optimal_lag)  # optimal_lag = 1
var_forecast = var_fitted.forecast(
    train.values[-optimal_lag:], steps=len(test))
\end{lstlisting}
\end{frame}

% ============================================================
\section{So sánh ARIMA vs VAR}
% ============================================================

\begin{frame}{So sánh ARIMA vs VAR -- Chỉ số đánh giá}
\begin{table}
\centering
\footnotesize
\begin{tabular}{llrrr}
\toprule
\textbf{Biến} & \textbf{Chỉ số} & \textbf{ARIMA} & \textbf{VAR} & \textbf{Tốt hơn} \\
\midrule
\multirow{4}{*}{CO(GT)}
    & MSE  & 0.7247  & 0.8769  & \textcolor{hkustblue}{ARIMA} \\
    & RMSE & 0.8513  & 0.9364  & \textcolor{hkustblue}{ARIMA} \\
    & MAE  & 0.6944  & 0.7539  & \textcolor{hkustblue}{ARIMA} \\
    & MAPE & 55.79\% & 63.09\% & \textcolor{hkustblue}{ARIMA} \\
\midrule
\multirow{4}{*}{C6H6(GT)}
    & MSE  & 24.5722 & 29.5813 & \textcolor{hkustblue}{ARIMA} \\
    & RMSE & 4.9570  & 5.4389  & \textcolor{hkustblue}{ARIMA} \\
    & MAE  & 4.1065  & 4.4438  & \textcolor{hkustblue}{ARIMA} \\
    & MAPE & 89.60\% & 97.72\% & \textcolor{hkustblue}{ARIMA} \\
\midrule
\multirow{4}{*}{NO2(GT)}
    & MSE  & 1662.15 & 1546.52 & \textcolor{hkustred}{VAR} \\
    & RMSE & 40.7695 & 39.3258 & \textcolor{hkustred}{VAR} \\
    & MAE  & 34.1071 & 33.1683 & \textcolor{hkustred}{VAR} \\
    & MAPE & 27.58\% & 21.96\% & \textcolor{hkustred}{VAR} \\
\midrule
\multirow{4}{*}{Nhiệt độ (T)}
    & MSE  & 28.2882 & 35.1688 & \textcolor{hkustblue}{ARIMA} \\
    & RMSE & 5.3187  & 5.9303  & \textcolor{hkustblue}{ARIMA} \\
    & MAE  & 4.2237  & 4.7197  & \textcolor{hkustblue}{ARIMA} \\
    & MAPE & 49.15\% & 83.97\% & \textcolor{hkustblue}{ARIMA} \\
\bottomrule
\end{tabular}
\end{table}
\end{frame}

% ============================================================
\section{Tổng kết}
% ============================================================

\begin{frame}{Tổng kết -- Trung bình trên tất cả biến}
\begin{table}
\centering
\begin{tabular}{lrrr}
\toprule
\textbf{Chỉ số} & \textbf{ARIMA (TB)} & \textbf{VAR (TB)} & \textbf{Tốt hơn} \\
\midrule
MSE  & 428.93 & 403.04 & \textcolor{hkustred}{\textbf{VAR}} \\
RMSE & 12.97  & 12.91  & \textcolor{hkustred}{\textbf{VAR}} \\
MAE  & 10.78  & 10.77  & \textcolor{hkustred}{\textbf{VAR}} \\
MAPE & 55.53\% & 66.69\% & \textcolor{hkustblue}{\textbf{ARIMA}} \\
\bottomrule
\end{tabular}
\end{table}

\vspace{0.5cm}
\begin{itemize}
    \item ARIMA thắng \textbf{1/4} chỉ số
    \item VAR thắng \textbf{3/4} chỉ số
\end{itemize}

\vspace{0.3cm}
\begin{block}{Kết luận}
$\Rightarrow$ Mô hình \textbf{VAR} hoạt động tốt hơn tổng thể trên bộ dữ liệu này.
\end{block}
\end{frame}

% ============================================================
\section{Những điểm chính}
% ============================================================

\begin{frame}{Những điểm chính}
\begin{table}
\centering
\begin{tabular}{lcc}
\toprule
 & \textbf{ARIMA} & \textbf{VAR} \\
\midrule
\textbf{Kiểu mô hình} & Đơn biến & Đa biến \\
\textbf{Tương quan chéo} & Không & Có \\
\textbf{Tham số} & $(p, d, q)$ cho mỗi chuỗi & Bậc trễ cho tất cả chuỗi \\
\textbf{Phù hợp nhất cho} & Chuỗi độc lập & Chuỗi có tương quan \\
\bottomrule
\end{tabular}
\end{table}

\vspace{0.5cm}
\begin{itemize}
    \item \textbf{ARIMA}: Mỗi biến được mô hình hóa độc lập; tốt khi chuỗi có động lực riêng mạnh (CO, C6H6, T).
    \item \textbf{VAR}: Nắm bắt cách các biến ảnh hưởng lẫn nhau (ví dụ: CO và NO$_2$ đều từ khí thải giao thông); hưởng lợi từ các mẫu chung.
\end{itemize}
\end{frame}

% --- Thank you slide ---
\begin{frame}
\begin{center}
{\Large Cảm ơn cô đã lắng nghe!}
\vspace{1cm}

{\normalsize Vũ Công Tuấn Dương -- B22DCKH024}

{\small Khoa CNTT1 -- PTIT}
\end{center}
\end{frame}

\end{document}
